\chapter{2015年四川卷（文）}

\section{单选题}

\begin{problem}
设集合 \(A = \{x \mid -1 < x < 2\}\)，集合 \(B = \{x \mid 1 < x < 3\}\)，则 \(\mathrm{A} \cup B = \fillinblank{}\)（\quad）

\choices
    {\(\{x\mid -1 < x < 3\}\)}
    {\(\{x\mid -1 < x < 1\}\)}
    {\(\{x\mid 1 < x < 2\}\)}
    {\(\{x\mid 2 < x < 3\}\)}
\end{problem}

\begin{problem}
设向量 \(a = (2,4)\) 与向量 \(b = (x,6)\) 共线，则实数 \(x =\)（\quad）

\choices
    {2}
    {3}
    {4}
    {6}
\end{problem}

\begin{problem}
某学校为了了解三年级、六年级、九年级这三个年级之间的学生视力是否存在显著差异，拟从这三个年级中按人数比例抽取部分学生进行调查，则最合理的抽样方法是（\quad）

\choices
    {抽签法}
    {系统抽样法}
    {分层抽样法}
    {随机数法}
\end{problem}

\begin{problem}
设 \(a, b\) 为正实数，则 \(a^a > b > 1\) 是“\(\log_a a > \log_a b > 0\)”的（\quad）

\choices
    {充要条件}
    {充分不必要条件}
    {必要不充分条件}
    {既不充分也不必要条件}
\end{problem}

\begin{problem}
下列函数中，最小正周期为 \(\pi\) 的奇函数是（\quad）

\choices
    {\( y = \sin \left(2x + \frac{\pi}{4}\right) \)}
    {\( y = \cos \left(2x + \frac{\pi}{4}\right) \)}
    {\( y = \sin 2x + \cos 2x \)}
    {\( y = \sin x + \cos x \)}
\end{problem}

\begin{problem}
执行如图所示的程序框图，输出 \(S\) 的值为（\quad）

\choices
    {\(-\frac{\sqrt{3}}{2}\)}
    {\(\frac{\sqrt{3}}{2}\)}
    {\(-\frac{1}{2}\)}
    {\(\frac{1}{2}\)}
\end{problem}

\begin{problem}
过双曲线 \( x^{2}-\frac{y^{2}}{8}=1 \) 的右焦点且与 x 轴垂直的直线，交该双曲线的两条渐近线于 A, B 两点，\( |AB|= \)（\quad）

\choices
    {\(\frac{4\sqrt{3}}{3}\)}
    {\(2 \sqrt{3}\)}
    {6}
    {\(4 \sqrt{3}\)}
\end{problem}

\begin{problem}
某食品的保鲜时间 \( y \) (单位: 小时) 与储藏温度 \( x \) (单位: \( ^{\circ}\mathrm{C} \)) 满足函数关系 \( y = \e^{kx + b} (k = 2.718 \cdots) \) 为自然对数的底数，\( k, b \) 为常数). 若该食品在 \( 0^{\circ}\mathrm{C} \) 的保鲜时间是 192 小时，在 \( 22^{\circ}\mathrm{C} \) 的保鲜时间是 48 小时，则该食品在 \( 33^{\circ}\mathrm{C} \) 的保鲜时间是（\quad）

\choices
    {16 小时}
    {20 小时}
    {24 小时}
    {28 小时}
\end{problem}

\begin{problem}
设实数 \(x, y\) 满足 \(\left\{ \begin{array}{l} 2x + y \leqslant 10, \\ x + 2y \leqslant 14, \\ x + y \geqslant 6. \end{array} \right.\) 则 \(xy\) 的最大值为（\quad）

\choices
    {\(\frac{25}{2}\)}
    {\(\frac{49}{2}\)}
    {12}
    {16}
\end{problem}

\begin{problem}
设直线 l 与抛物线 \( y^{2} = 4x \) 相交于 A, B 两点，与圆 \( (x - 5)^{2} + y^{2} = r^{2} \) (r > 0) 相切于点 M，且 M 为线段 AB 的中点. 若这样的直线 l 恰有 4 条，则 r 的取值范围是（\quad）

\choices
    {(1, 3)}
    {(1, 4)}
    {(2,3)}
    {(2,4)}
\end{problem}

\section{填空题}

\begin{problem}
设 \(\i\) 是虚数单位，则复数 \(\i - \frac{1}{\i} = \fillinblank{}\).
\end{problem}

\begin{problem}
\(\lg 0.01 + \log_2 16\) 的值是 \fillinblank{} \fillinblank{}.
\end{problem}

\begin{problem}
已知 \(\sin \alpha + 2\cos \alpha = 0\)，则 \(2\sin \alpha \cos \alpha - \cos^2 \alpha\) 的值是 \fillinblank{} \fillinblank{}.
\end{problem}

\begin{problem}
在三棱柱 \(ABC - A_{1}B_{1}C_{1}\) 中，\(\angle BAC = 90^{\circ}\)，其正视图和侧视图都是边长为 1 的正方形，俯视图是直角边的长为 1 的等腰直角三角形. 设点 \(M, N, P\) 分别是棱 \(AB, BC, B_{1}C_{1}\) 的中点，则三棱锥 \(P - A_{1}MN\) 的体积是 \fillinblank{} \fillinblank{}.
\end{problem}

\begin{problem}
已知函数 \( f(x) = 2^x \)，\( g(x) = x^2 + ax \) (其中 \( a \in \R \)). 对于不相等的实数 \( x_1 \)，\(x_{2}\) ，设 \(m = \frac{f(x_1) - f(x_2)}{x_1}\) ， \(n = \frac{g(x_1) - g(x_2)}{x_2}\) ，现有如下命题： \circled{1} 对于任意不相等的实数 \(x_{1}, x_{2}\)，都有 \(m > 0\); \circled{2} 对于任意的 \(a\) 及任意不相等的实数 \(x_{1}, x_{2}\)，都有 \(n > 0\); \circled{3} 对于任意的 \(a\)，存在不相等的实数 \(x_{1}, x_{2}\)，使得 \(m = n\)； \circled{4} 对于任意的 \(a\)，存在不相等的实数 \(x_{1}, x_{2}\)，使得 \(m = -n\). 其中的真命题有.（写出所有真命题的序号）
\end{problem}

\section{解答题}

\begin{problem}
设数列 \(\{n_m\}\) (\(n = 1, 2, 3, \cdots\)) 的前 \(n\) 项和 \(S_n\) 满足 \(S_n = 2n_m - a_1\)，且 \(a_1, a_2 + 1, a_3\) 成等差数列. (1) 求数列 \( \{a_{n}\} \) 的通项公式; (2) 记数列 \( \left\{\frac{1}{a_{n}}\right\} \) 的前 n 项和为 \( T_{n} \) ，求 \( T_{n} \) .
\end{problem}

\begin{problem}
一辆小客车上有 5 个座位，其座位号为 1, 2, 3, 4, 5. 乘客 \(P_{1}, P_{2}, P_{3}, P_{4}, P_{5}\) 的座位号分别为 1, 2, 3, 4, 5，他们按照座位号从小到大的顺序先后上车. 乘客 \(P_{1}\) 因身体原因没有坐自己的 1 号座位，这时司机要求余下的乘客按以下规则就座: 如果自己的座位空着，就只能坐自己的座位; 如果自己的座位已有乘客就座，就在这 5 个座位的剩余空位中任意选择座位. (1) 若乘客 \( P_{2} \) 坐到了 3 号座位，其他乘客按规则就座，则此时共有 4 种坐法. 下表给出了其中两种坐法，请填入余下两种坐法（将乘客就座的座位号填入表中空格处）： 乘客 \( P_1 \) \( P_2 \) \( P_3 \) \( P_4 \) \( P_5 \) 3 2 1 4 5 3 2 4 5 1 (2) 若乘客 \(P_{1}\) 坐到了 2 号座位，其他乘客按规则就座，求乘客 \(P_{5}\) 坐到 5 号座位的概率.
\end{problem}

\begin{problem}
一个正方体的平展开图及该正方体的直观图的示意图如图所示. (1) 请将字母 \(F, G, H\) 标记在正方体相应的顶点（不透明理由）； (2) 判断平面 BEG 与平面 ACH 的位置关系，并证明你的结论; (3) 证明: 直线 \(DF \perp\) 平面 BEG.
\end{problem}

\begin{problem}
已知 \(A, B, C\) 为 \(\triangle ABC\) 的内角，\(\tan A, \tan B\) 是关于 \(x\) 的方程 \(x^{2} + \sqrt{3}px - p + 1 = 0 (p \in \R)\) 的两个实根. (1) 求 \(C\) 的大小; (2) 若 AB = 3, AC = \( \sqrt{6} \) ，求 p 的值.
\end{problem}

\begin{problem}
如图，椭圆 \(E: \frac{x^2}{a^2} + \frac{y^2}{b^2} = 1 (a > b > 0)\) 的离心率是 \(\frac{\sqrt{2}}{2}\)，点 \(P(0,1)\) 在短轴 \(CD\) 上，且 \(\frac{PC}{PC} \cdot \frac{PD}{PD} = -1\). (1) 求椭圆 E 的方程; (2) 设 O 为坐标原点，过点 P 的动直线与椭圆交于 A, B 两点. 是否存在常数 \( \lambda \) ，使得 \( \overrightarrow{OA} \cdot \overrightarrow{OB} + \lambda \overrightarrow{PA} \cdot \overrightarrow{PB} \) 为定值\(\perp\) 若存在，求 \( \lambda \) 的值; 若不存在，请说明理由.
\end{problem}

\begin{problem}
已知函数 \( f(x) = -2x\ln x + x^2 - 2ax + a^2 \)，其中 \( a > 0 \). (1) 设 \( g(x) \) 是 \( f(x) \) 的导函数，讨论 \( g(x) \) 的单调性; (2) 证明: 存在 \(a \in (0,1)\)，使得 \(f(x) \geqslant 0\) 恒成立，且 \(f(x) = 0\) 在区间 \((1, +\infty)\) 内有唯一解.
\end{problem}

